% ============================================================
\chapter{Standard Distributions and Properties}
\label{ch:distributions}
% ============================================================

Probability theory has produced, over the centuries, a gallery of remarkable distributions---each born from a concrete problem, and each endowed with specific properties that make it irreplaceable in its domain. The Gaussian normal, the exponential of waiting times, the Poisson of rare events, the Beta of Bayesian statistics, the Gamma of queuing theory: this chapter gathers them in an organised panorama, highlighting their mutual relationships (sums, conditioning, limits) and the properties that characterise them.

\begin{intuition}
This chapter presents an organised catalogue of the most common probability
distributions, both discrete and continuous, together with their properties,
mutual relationships, and domains of application.
\end{intuition}

% ------------------------------------------------------------
\section{Discrete Distributions}
% ------------------------------------------------------------

\subsection{Negative binomial distribution}

\begin{definition}[Negative binomial]
$X \sim \mathrm{NB}(r, p)$: number of trials to obtain the $r$-th success.
\[
    \PP(X = k) = \binom{k-1}{r-1} p^r(1-p)^{k-r}, \quad k = r, r+1, \ldots
\]
$\E[X] = r/p$, $\Var(X) = r(1-p)/p^2$. For $r=1$ this reduces to the geometric
distribution.
\end{definition}

\subsection{Discrete uniform distribution}

\begin{definition}[Discrete uniform]
$X \sim \mathcal{U}\{a, a+1, \ldots, b\}$:
\[
    \PP(X = k) = \frac{1}{b-a+1}, \quad k = a, a+1, \ldots, b.
\]
$\E[X] = (a+b)/2$, $\Var(X) = ((b-a+1)^2-1)/12$.
\end{definition}

\subsection{Poisson distribution --- further properties}

\begin{proposition}[Stability under summation]
If $X_1 \sim \mathrm{Poisson}(\lambda_1)$ and $X_2 \sim \mathrm{Poisson}(\lambda_2)$
are independent, then $X_1 + X_2 \sim \mathrm{Poisson}(\lambda_1+\lambda_2)$.
\end{proposition}

\begin{proposition}[Poisson splitting]
If $X \sim \mathrm{Poisson}(\lambda)$ and each event is independently classified
as type~1 with probability $p$ or type~2 with probability $1-p$, then the counts
$N_1, N_2$ are independent with $N_1 \sim \mathrm{Poisson}(\lambda p)$ and
$N_2 \sim \mathrm{Poisson}(\lambda(1-p))$.
\end{proposition}

% ------------------------------------------------------------
\section{Continuous Distributions}
% ------------------------------------------------------------

\subsection{Log-normal distribution}

\begin{definition}[Log-normal]
$X \sim \mathrm{LogN}(\mu, \sigma^2)$ if $\ln X \sim \mathcal{N}(\mu, \sigma^2)$.
\[
    f_X(x) = \frac{1}{x\sigma\sqrt{2\pi}}
    \exp\!\left(-\frac{(\ln x - \mu)^2}{2\sigma^2}\right), \quad x > 0.
\]
$\E[X] = e^{\mu+\sigma^2/2}$, $\Var(X) = e^{2\mu+\sigma^2}(e^{\sigma^2}-1)$.
\end{definition}

\subsection{Cauchy distribution}

\begin{definition}[Cauchy]
$X \sim \mathrm{Cauchy}(x_0, \gamma)$:
\[
    f_X(x) = \frac{1}{\pi\gamma\left(1 + \left(\frac{x-x_0}{\gamma}\right)^2\right)}.
\]
The expectation and variance do not exist.
\end{definition}

\begin{warning}[title=\textbf{Cauchy has no moments}]
The Cauchy distribution shows that not every random variable has an expectation.
The sample mean of i.i.d.\ Cauchy variables does not converge: the law of large
numbers does not apply.
\end{warning}

\subsection{Weibull distribution}

\begin{definition}[Weibull]
$X \sim \mathrm{Weibull}(k, \lambda)$ with $k, \lambda > 0$:
\[
    f_X(x) = \frac{k}{\lambda}\left(\frac{x}{\lambda}\right)^{k-1}
    e^{-(x/\lambda)^k}\,\mathbf{1}_{(0,+\infty)}(x).
\]
For $k=1$ this is the exponential distribution with parameter $1/\lambda$.
\end{definition}

\subsection{Chi-squared distribution}

\begin{definition}[Chi-squared]
If $Z_1, \ldots, Z_n$ are i.i.d.\ $\mathcal{N}(0,1)$, then
$\chi^2_n = Z_1^2 + \cdots + Z_n^2 \sim \Gamma(n/2, 1/2)$ is the
\textbf{chi-squared distribution with $n$ degrees of freedom}.
$\E[\chi^2_n] = n$, $\Var(\chi^2_n) = 2n$.
\end{definition}

\subsection{Student's $t$-distribution}

\begin{definition}[Student's $t$]
If $Z \sim \mathcal{N}(0,1)$ and $V \sim \chi^2_n$ are independent, then
$T = Z/\sqrt{V/n}$ follows the \textbf{Student $t$-distribution with $n$ degrees
of freedom}. $\E[T] = 0$ (for $n > 1$), $\Var(T) = n/(n-2)$ (for $n > 2$).
For $n=1$, one recovers the standard Cauchy distribution.
\end{definition}

\subsection{Fisher's $F$-distribution}

\begin{definition}[Fisher's $F$]
If $U \sim \chi^2_m$ and $V \sim \chi^2_n$ are independent, then
$F = (U/m)/(V/n)$ follows the \textbf{$F$-distribution with $(m, n)$ degrees
of freedom}.
\end{definition}

% ------------------------------------------------------------
\section{Relationships Between Distributions}
% ------------------------------------------------------------

\begin{keyformulas}[title=\textbf{Map of distribution relationships}]
\begin{itemize}
    \item $\mathrm{Bernoulli}(p) = \mathrm{Bin}(1, p)$.
    \item $\mathrm{Geom}(p) = \mathrm{NB}(1, p)$.
    \item $\mathrm{Exp}(\lambda) = \Gamma(1, \lambda)$.
    \item $\chi^2(n) = \Gamma(n/2, 1/2)$.
    \item $\mathrm{Bin}(n,p) \xrightarrow{n\to\infty,\,np\to\lambda} \mathrm{Poisson}(\lambda)$.
    \item $\mathrm{Bin}(n,p) \xrightarrow{n\to\infty} \mathcal{N}(np, np(1-p))$ \quad (CLT).
    \item $\mathrm{Poisson}(\lambda) \xrightarrow{\lambda\to\infty} \mathcal{N}(\lambda, \lambda)$.
    \item $t_n \xrightarrow{n\to\infty} \mathcal{N}(0,1)$.
    \item $\Gamma(n,\lambda)$ = distribution of the sum of $n$ i.i.d.\ $\mathrm{Exp}(\lambda)$ variables.
\end{itemize}
\end{keyformulas}

% ------------------------------------------------------------
\section{General Summary Table}
% ------------------------------------------------------------

\begin{keyformulas}[title=\textbf{Continuous distributions summary}]
\begin{center}
\renewcommand{\arraystretch}{1.4}
\begin{tabular}{lcccc}
    \toprule
    \textbf{Distribution} & \textbf{Density} & \textbf{Support} & $\E[X]$ & $\Var(X)$ \\
    \midrule
    $\mathcal{U}([a,b])$ & $\frac{1}{b-a}$ & $[a,b]$ & $\frac{a+b}{2}$ & $\frac{(b-a)^2}{12}$ \\[4pt]
    $\mathrm{Exp}(\lambda)$ & $\lambda e^{-\lambda x}$ & $[0,\infty)$ & $\frac{1}{\lambda}$ & $\frac{1}{\lambda^2}$ \\[4pt]
    $\mathcal{N}(\mu,\sigma^2)$ & $\frac{1}{\sigma\sqrt{2\pi}}e^{-(x-\mu)^2/(2\sigma^2)}$ & $\R$ & $\mu$ & $\sigma^2$ \\[4pt]
    $\Gamma(\alpha,\lambda)$ & $\frac{\lambda^\alpha}{\Gamma(\alpha)}x^{\alpha-1}e^{-\lambda x}$ & $(0,\infty)$ & $\frac{\alpha}{\lambda}$ & $\frac{\alpha}{\lambda^2}$ \\[4pt]
    $\mathrm{Beta}(\alpha,\beta)$ & $\frac{x^{\alpha-1}(1-x)^{\beta-1}}{B(\alpha,\beta)}$ & $(0,1)$ & $\frac{\alpha}{\alpha+\beta}$ & $\frac{\alpha\beta}{(\alpha+\beta)^2(\alpha+\beta+1)}$ \\
    \bottomrule
\end{tabular}
\end{center}
\end{keyformulas}

% ------------------------------------------------------------
\section{Identification Methods}
% ------------------------------------------------------------

\begin{remark}
To identify a distribution, proceed as follows:
\begin{enumerate}
    \item Identify the \textbf{support}: $\N$, $\{0,1,\ldots,n\}$, $(0,\infty)$,
          $\R$, $(0,1)$, etc.
    \item Recognise the \textbf{functional form}: decaying exponential, symmetric
          bell curve, power law, etc.
    \item Check the \textbf{parameters} by computing the mean and variance.
\end{enumerate}
\end{remark}

% ------------------------------------------------------------
\section{Exercises}
% ------------------------------------------------------------

\begin{exercise}
Let $X \sim \Gamma(n, \lambda)$ with $n \in \N^*$. Show that $X$ has the same
distribution as the sum of $n$ i.i.d.\ $\mathrm{Exp}(\lambda)$ variables.
\end{exercise}

\begin{exercise}
Show that if $X \sim \mathcal{N}(0,1)$, then $X^2 \sim \chi^2(1)$.
\end{exercise}

\begin{exercise}
Let $X \sim \mathrm{Poisson}(\lambda)$ with $\lambda = 100$. Using the normal
approximation, estimate $\PP(90 \leq X \leq 110)$.
\end{exercise}

\begin{exercise}
Let $U \sim \mathcal{U}([0,1])$. Find the distribution of $X = -\ln U$.
\end{exercise}

\begin{exercise}
Let $X \sim \mathrm{Beta}(1,1)$. Show that $X \sim \mathcal{U}([0,1])$.
\end{exercise}

\begin{exercise}
Customers arrive at a shop according to a Poisson process with rate
$\lambda = 12$ per hour. Each customer independently makes a purchase with
probability $p = 0.3$. What is the distribution of the number of purchasing
customers per hour?
\end{exercise}

\begin{exercise}
Show that if $X \sim t_n$ (Student with $n$ df), then $X^2 \sim F(1, n)$ (Fisher).
\end{exercise}

\begin{exercise}
Let $X \sim \mathrm{LogN}(0, 1)$. Compute $\E[X]$, $\Var(X)$, and the median of $X$.
\end{exercise}
